% Copyright (c) 2026 Geovana Neves. Licensed under CC BY 4.0 (https://creativecommons.org/licenses/by/4.0/). See LICENSE-AND-AUTHORSHIP.md.
%
\section{The sections table}

\begin{frame}{One export contains every declared distribution}
\begin{itemize}
  \item \code{sections/<point>\_sections.csv} is one solver export.
  \item It holds every distribution the pproc declares: the wing in
        \code{XZ}, each blade in its own frame.
  \item The distributions follow one another with no marker between them.
  \item Leading columns identify which distribution each row belongs to.
\end{itemize}
\end{frame}

\begin{frame}{The sections table: step and distribution identity}
\relax
{\ftstablesize
\begin{tabular}{@{}p{0.17\linewidth}p{0.76\linewidth}@{}}
\toprule
\textbf{Column} & \textbf{Meaning} \\
\midrule
\code{STEP} & Unsteady: the run's last time step, from the run record.
The export is written when the march ends; its own header counts inner
iterations, not time steps. Steady: the solver iteration the export states.
One name across every table \\
\code{FAMILY} & Geometry families of the row's distribution, joined by \code{+} \\
\code{PLANE} & Cutting plane of that distribution \\
\bottomrule
\end{tabular}}
\end{frame}

\begin{frame}{The sections table: rotor identity}
\relax
{\ftstablesize
\begin{tabular}{@{}p{0.17\linewidth}p{0.76\linewidth}@{}}
\toprule
\textbf{Column} & \textbf{Meaning, after \code{PLANE}} \\
\midrule
\code{ROTOR} & Rotor whose blades those families are;
\code{NA} for a surface no rotor owns \\
\code{AZIMUTH} & Where blade one of that rotor is at \code{STEP}, in degrees,
wrapped to one turn; \code{NA} without a rotor \\
\bottomrule
\end{tabular}}
\end{frame}

\begin{frame}{Each rotor has its own azimuth clock}
\small
\[
\mathrm{AZIMUTH}=\left(
  \mathrm{blade1.azimuth\_deg}
  +\mathrm{sense}\,\mathrm{STEP}\,
   \frac{360}{\mathrm{steps\_per\_revolution}}
\right)\bmod 360.
\]
\begin{itemize}
  \item The datum, the sense of rotation and the steps per revolution all
        come from \textbf{that rotor}.
  \item Sense is the sign of the rotor's speed.
  \item Two rotors at two speeds have two azimuths at the same step.
        A wing has none.
\end{itemize}
\end{frame}

\begin{frame}{Distribution identity comes from the run record}
\begin{itemize}
  \item The run records which distribution is which. The script states
        surfaces by index; post cannot recover their names from that alone.
  \item An older record writes \code{NA} in all four identity
        fields: \code{FAMILY}, \code{PLANE}, \code{ROTOR}, \code{AZIMUTH}.
  \item The same rule applies when the recorded blocks do not add up to
        the number of exported rows.
\end{itemize}
\end{frame}

\begin{frame}{The sections table is one instant}
\begin{itemize}
  \item At an unsteady point, the distribution is at \code{STEP}.
        It is not an average over the window.
  \item Its \code{products.json} entry states \code{"kind": "instant"}.
  \item The history is \code{series/<point>\_sections\_series.csv},
        with the same identity on every row.
\end{itemize}
\end{frame}

\begin{frame}{One loads file and one Cp file per distribution}
\begin{itemize}\small
  \item Each \code{[[sections.distributions]]} entry gets
        \code{sections/<point>\_sloads\_<name>.csv} and
        \code{sections/<point>\_cp\_<name>.csv} when its export is available.
  \item The name preserves the alias or joins the declared families with
        \code{-}. Filename-invalid characters become \code{\_}; collisions
        receive the entry's 1-based position suffix.
  \item All planes and expanded blade blocks of that entry share its file.
        Cp now has a table, rather than only a manifest listing.
  \item With per-step exports enabled, every exported step appears in ascending
        \code{STEP} order. Without them, the file holds the end-of-run export.
  \item The combined sections table and combined sections series remain available.
\end{itemize}
\end{frame}

\begin{frame}{Distribution columns and recorded identity}
\begin{itemize}\small
  \item Rows start with \code{STEP, time\_s, FAMILY, PLANE, ROTOR, AZIMUTH},
        then the condition block and export columns. Cp also carries
        \code{SECTION}, the export's 1-based cross-section index.
  \item Missing values are \code{NA}. The manifest records the distribution,
        original families and \code{steps\_tabled}.
  \item A new run records each block's distribution entry. A record can
        be split only when its recorded layout and pproc identify the blocks
        unambiguously by families, plane, frame and count.
  \item Missing or ambiguous identity, mismatched counts, malformed exports
        and unavailable steps are named skips in \code{products.json}.
        Editing today's pproc cannot supply missing recorded identity.
\end{itemize}
\end{frame}
